\documentclass{article}
\usepackage[utf8]{inputenc}
\usepackage{amsmath}
\usepackage{amsthm}
\usepackage{amssymb,amsfonts}
\usepackage{mathrsfs}
\usepackage[hidelinks]{hyperref}
\usepackage[capitalize]{cleveref}
\usepackage{stmaryrd}

\newtheorem{theorem}{Theorem}[section]
\newtheorem{proposition}[theorem]{Proposition}
\newtheorem{claim}[theorem]{Claim}
\newtheorem{corollary}[theorem]{Corollary}
\newtheorem{lemma}[theorem]{Lemma}
\theoremstyle{definition}
\newtheorem{definition}[theorem]{Definition}
\newtheorem{convention}[theorem]{Convention}
\newtheorem{remark}[theorem]{Remark}

\newcommand{\Rbar}{\overline{\mathbb{R}}}
\newcommand{\R}{\mathbb{R}}
\newcommand{\rk}{\operatorname{rk}}
\newcommand{\cl}{\operatorname{cl}}
\newcommand{\Trop}{\operatorname{Trop}}
\newcommand{\op}{\operatorname{op}}
\newcommand{\into}{\hookrightarrow}
\newcommand{\onto}{\twoheadrightarrow}
\newcommand{\Dr}{\mathbf{Dr}}
\newcommand{\bV}{\mathbf V}
\newcommand{\cLL}{\mathcal{L}}
\newcommand{\cW}{\mathcal{W}}
\newcommand{\PT}[1]{\mathbf{P}\big(\Rbar^{#1}\big)}

\title{First proof: Duality for tropical linear subspaces}
\author{}
\date{}

\begin{document}
\maketitle

\noindent
Let $M$ be a matroid on the ground set $E=[n]$, of rank $r$, and let $\Trop(M)$ be the
corresponding tropical linear space, i.e.\ the closure of the Bergman fan of $M$ in tropical
projective space $\PT{n}$. The \emph{relative Dressian} of $M$ is the set
\[
  \Dr(M):=\{\text{tropical linear subspaces of }\Trop(M)\},
\]
partially ordered by inclusion of tropical linear spaces. Suppose $M$ admits an order-reversing
involution $F\mapsto F^{\perp}$ on its lattice of flats $\cLL(M)$. There is a natural inclusion of
posets $\cLL(M)^{\op}\into\Dr(M)$ given by $F\mapsto\Trop(M/F\oplus U_{0,F})$.

\medskip
\noindent\textbf{Theorem.} \emph{There is an order-reversing involution $\Phi$ on $\Dr(M)$ that
extends the involution $F\mapsto F^{\perp}$; that is, $\Phi$ restricts to $F\mapsto F^{\perp}$ on the
image of $\cLL(M)^{\op}\into\Dr(M)$.}

\medskip
\noindent\textbf{Status of this document.} This file records a \emph{partial, honest} treatment of
the Theorem. The construction of the candidate involution $\Phi$ is given explicitly, and several of
its required properties are proved in full; one decisive property --- the well-definedness of $\Phi$,
i.e.\ that the property of being a valuated quotient of the trivial valuation is invariant under the
operation that $\perp$ induces on flat-coordinates --- is reduced to a single, precisely isolated
claim (\Cref{conj:wd-invariance}) which we are unable to prove in full generality and which appears
to be the genuine technical content of the announced reference. We have \emph{verified the whole
Theorem computationally} on the smallest nontrivial modular matroid carrying a polarity (the Fano
plane $F_7$), as well as the rank/lattice identities and the constancy and extension statements; see
\Cref{rmk:honest,rmk:numeric}. We have also \emph{disproved by explicit counterexample} the
``essential relations'' reduction on which a naive proof (and an earlier draft) rested; see
\Cref{rmk:false-reduction}. Thus the present write-up is faithful: it proves what is provable by the
methods at hand, and it flags --- rather than papers over --- the remaining gap.

\medskip

\noindent\textbf{Outline.} We fix conventions and the two standard foundational facts
(\Cref{ssec:found}); prove the rank identity and lattice anti-automorphism forced by $\perp$
(\Cref{lem:rank}); show that every element of $\Dr(M)$ has a canonical \emph{flat-coordinate}
description, the underlying valuated quotient being constant on the bases lying in a fixed flat
(\Cref{prop:constancy}); define the candidate $\Phi$ explicitly as precomposition of flat-coordinates
with $\perp$ (\Cref{def:Phi}); isolate the well-definedness statement (\Cref{conj:wd-invariance})
and prove that, \emph{granting} it, $\Phi$ is an involution (\Cref{prop:inv}); prove the extension
property $\Phi(L_F)=L_{F^{\perp}}$ unconditionally (\Cref{prop:ext}); and record the honest status
and the computational evidence (\Cref{rmk:honest,rmk:false-reduction,rmk:numeric}).

\section{Conventions and foundational facts}\label{ssec:found}

We use the min-plus conventions on $\Rbar=\R\cup\{\infty\}$. A \emph{valuated matroid} of rank $k$
on $E$ is a function $\mu\colon\binom{E}{k}\to\Rbar$, not identically $\infty$, satisfying the
tropical Pl\"ucker relations: for every $S\in\binom{E}{k-1}$ and every $T\in\binom{E}{k+1}$,
\begin{equation}\label{eqn:plucker}
  \min_{a\in T\setminus S}\bigl\{\mu(S+a)+\mu(T-a)\bigr\}\ \text{vanishes tropically,}
\end{equation}
i.e.\ either all summands are $\infty$, or the minimum is attained for at least two
$a\in T\setminus S$. Here $S+a:=S\cup\{a\}$ and $T-a:=T\setminus\{a\}$. Two valuated matroids of the
same rank are \emph{equivalent} if they differ by a global additive constant; the support
$\{B:\mu(B)<\infty\}$ is the set of bases of the \emph{underlying matroid} $\underline\mu$. The
\emph{trivial valuation} $\underline M$ of $M$ is the rank-$r$ valuated matroid with $\underline
M(B)=0$ if $B$ is a basis of $M$ and $\underline M(B)=\infty$ otherwise; thus
$\underline{\underline M}=M$.

\begin{definition}\label{def:quotient}
Let $\nu,\mu$ be valuated matroids on $E$ of ranks $k\le m$. We say $\nu$ is a \emph{valuated
quotient} of $\mu$, written $\mu\onto\nu$, if for every $S\in\binom{E}{k-1}$ and every
$T\in\binom{E}{m+1}$,
\begin{equation}\label{eqn:incidence-plucker}
  \min_{a\in T\setminus S}\bigl\{\nu(S+a)+\mu(T-a)\bigr\}\ \text{vanishes tropically.}
\end{equation}
Here $\nu$ is evaluated on the $k$-set $S+a$ and $\mu$ on the $m$-set $T-a$; the index sets have the
correct cardinalities $|S+a|=k$ and $|T-a|=m$, since $|S|=k-1$ and $|T|=m+1$. When $m=k+1$ this is
the \emph{elementary} quotient relation; the relation \eqref{eqn:incidence-plucker} for general
$k\le m$ is equivalent to the existence of a chain of elementary quotients
$\mu=\mu_m\onto\mu_{m-1}\onto\cdots\onto\mu_k=\nu$ (Murota, Brandt--Eur--Speyer--Sturmfels).
\end{definition}

\begin{proposition}[Dressian as poset of valuated quotients]\label{fact:dressian}
Let $M$ be loopless of rank $r$. The assignment $\theta\mapsto\Trop\theta$ is an isomorphism of
posets
\[
  \bV(M)/\!\sim\ \xrightarrow{\ \cong\ }\ \Dr(M),
\]
where $\bV(M):=\bigsqcup_{0\le k\le r}\bV^{k}(M)$, $\bV^{k}(M)$ is the set of corank-$k$ valuated
quotients of $\underline M$ (i.e.\ valuated quotients of rank $r-k$), $\sim$ is equivalence up to a
global additive constant, the order on $\bV(M)/\!\sim$ is the quotient relation
($\theta_2\ge\theta_1\iff\theta_2\onto\theta_1$), and the order on $\Dr(M)$ is inclusion. In
particular $\theta_2\onto\theta_1\iff\Trop\theta_1\subseteq\Trop\theta_2$.
\end{proposition}

\begin{proof}
This is the relative form of the standard correspondence between tropical linear subspaces and
valuated matroids, and between inclusions and valuated quotients (Maclagan--Sturmfels,
\emph{Introduction to Tropical Geometry}, Ch.~4; Brandt--Eur--Sturmfels; Brandt--Speyer). We verify
its standing hypothesis that $M$ is loopless. A loop of $M$ lies in $\cl(\varnothing)$, the bottom
flat $\hat0$ of $\cLL(M)$. The order-reversing bijection $\perp$ sends the unique top
$\hat1=E$ to the unique bottom, so $\hat1^{\perp}=\hat0$ and $\hat0^{\perp}=\hat1$; in particular
$\cLL(M)$ has length $r$ from $\hat0$ to $\hat1$, so $\hat0$ has rank $0$ and contains no nonloop.
If $M$ had loops $L$, then $\Trop(M)$, $\cLL(M)$, $\perp$, and $\Dr(M)$ are unchanged upon deleting
$L$ (loops are common to all bases' complements and carry no coordinate in $\Trop M$); we may thus
assume $M$ loopless. The order-preservation in both directions is the cited quotient/containment
dictionary.
\end{proof}

\begin{proposition}[Valuated matroid duality]\label{fact:dual}
Let $\theta$ be a valuated matroid of rank $k$ on $E$, $|E|=n$. Define
$\theta^{*}\colon\binom{E}{n-k}\to\Rbar$ by $\theta^{*}(B):=\theta(E\setminus B)$. Then $\theta^{*}$
is a valuated matroid of rank $n-k$ with $\underline{\theta^{*}}=(\underline\theta)^{*}$ and
$(\theta^{*})^{*}=\theta$, and $\mu\onto\nu$ implies $\nu^{*}\onto\mu^{*}$.
\end{proposition}

\begin{proof}
Dress--Wenzel duality. The complementation $B\mapsto E\setminus B$ identifies $\binom{E}{n-k}$ with
$\binom{E}{k}$ and, for $a\in T\setminus S$, satisfies $E\setminus(S+a)=(E\setminus S)-a$ and
$E\setminus(T-a)=(E\setminus T)+a$; thus it carries the Pl\"ucker relation \eqref{eqn:plucker} of
$\theta$ at $(S,T)$ term-by-term to the Pl\"ucker relation of $\theta^{*}$ at $(E\setminus
T,E\setminus S)$, and similarly carries the incidence relation \eqref{eqn:incidence-plucker} for
$\mu\onto\nu$ to that for $\nu^{*}\onto\mu^{*}$. Double duality is immediate from $E\setminus
(E\setminus B)=B$.
\end{proof}

\begin{convention}\label{conv:perp}
Throughout, $\perp$ is a fixed inclusion-reversing bijection $\cLL(M)\to\cLL(M)$ with
$(F^{\perp})^{\perp}=F$. We do \emph{not} assume $M$ is a modular matroid in the strong sense (every
pair of flats a modular pair); we use only the existence of $\perp$ and its consequences derived
below.
\end{convention}

\section{The rank identity and the lattice anti-automorphism}

\begin{lemma}\label{lem:rank}
For every $F\in\cLL(M)$, $\rk_M(F)+\rk_M(F^{\perp})=r$. Moreover $\hat0^{\perp}=\hat1$ and
$\hat1^{\perp}=\hat0$, where $\hat0=\cl(\varnothing)$, $\hat1=E$. The map $\perp$ is an
anti-automorphism of the lattice $\cLL(M)$: $(F\vee G)^{\perp}=F^{\perp}\wedge G^{\perp}$ and
$(F\wedge G)^{\perp}=F^{\perp}\vee G^{\perp}$, and it restricts to a bijection
$\cLL_k(M)\to\cLL_{r-k}(M)$ for every $k$, where $\cLL_k(M)$ denotes the rank-$k$ flats.
\end{lemma}

\begin{proof}
$\cLL(M)$ is graded of rank $r$: every maximal chain $\hat0=G_0\lessdot\cdots\lessdot G_r=\hat1$ has
length $r$, and $\rk_M(G)$ equals the height of $G$. An order-reversing bijection sends covers to
covers ($A\lessdot B\iff B^{\perp}\lessdot A^{\perp}$), hence maximal chains to maximal chains in
reverse order. Applying $\perp$ to a maximal chain through $F$ at height $\rk_M(F)$ yields a maximal
chain in which $F^{\perp}$ sits at height $r-\rk_M(F)$, so $\rk_M(F^{\perp})=r-\rk_M(F)$; in
particular $\perp$ maps $\cLL_k(M)$ onto $\cLL_{r-k}(M)$. Taking $F=\hat0$ gives $\hat0^{\perp}=\hat1$,
and involutivity gives $\hat1^{\perp}=\hat0$. Finally, in any lattice an order-reversing bijection
exchanges joins and meets, because $F\vee G$ is the least upper bound of $\{F,G\}$, whose image under
the order-reversing $\perp$ is the greatest lower bound of $\{F^{\perp},G^{\perp}\}$, namely
$F^{\perp}\wedge G^{\perp}$; and dually.
\end{proof}

\section{Flat-coordinates: constancy on flats}

\begin{proposition}[Constancy on flats]\label{prop:constancy}
Let $\theta\in\bV^{\,r-k}(M)$ be a valuated quotient of $\underline M$ of rank $k$, with underlying
matroid $N:=\underline\theta$. If $B,B'$ are bases of $N$ with $\cl_M(B)=\cl_M(B')$, then
$\theta(B)=\theta(B')$. Consequently the function
\begin{equation}\label{eqn:hat-def}
  \widehat\theta\colon\cLL_k(M)\to\Rbar,\qquad
  \widehat\theta(F):=\begin{cases}\theta(B), & \text{$F=\cl_M(B)$ for a basis $B$ of $N$},\\
  \infty, & \text{no basis $B$ of $N$ has $\cl_M(B)=F$,}\end{cases}
\end{equation}
is well-defined up to the global additive constant of $\theta$, and $\theta(B)=\widehat\theta(\cl_M
B)$ for every basis $B$ of $N$. (If $B$ is a basis of $N$ then $B$ is independent in $M$, so
$\cl_M(B)\in\cLL_k(M)$.)
\end{proposition}

\begin{proof}
Since $M\onto N$ is a matroid quotient, $\rk_N\le\rk_M$ and every $N$-independent set is
$M$-independent; in particular a basis $B$ of $N$ (size $k$) is $M$-independent, so
$\cl_M(B)\in\cLL_k(M)$. This proves the parenthetical claim.

For the constancy, it suffices to treat $B'=B-x+y$ obtained from $B$ by a single symmetric exchange
($x\in B$, $y\notin B$, $B'$ a basis of $N$) with $\cl_M(B)=\cl_M(B')=:F$: the bases of $N$ that lie
in $F$ are exactly the bases of the restriction $N|_F$, whose basis-exchange graph is connected, and
along every exchange edge inside $F$ both endpoints have $M$-closure $F$. So assume $B'=B-x+y$ with
$\cl_M(B)=\cl_M(B')=F$; thus $y\in\cl_M(B)\setminus B$, i.e.\ $B+y$ has $M$-rank $k$ (it lies in $F$).

Apply the quotient relation \eqref{eqn:incidence-plucker} for $\underline M\onto\theta$ (so $m=r$,
the relation of $\theta$ with the trivial valuation) with
\[
  S:=B-x\quad(|S|=k-1),\qquad T:=(B-x)\cup\{x,y,z_1,\dots,z_{r-k}\}\quad(|T|=r+1),
\]
where $z_1,\dots,z_{r-k}\in E$ are chosen so that $B\cup\{z_1,\dots,z_{r-k}\}$ is a basis of $M$.
Such $z_i$ exist because $\rk_M(B)=k$ and $M$ has rank $r$ (extend the independent set $B$ to an
$M$-basis). Note $z_i\notin F$ for all $i$, since adding any $z_i$ raises $M$-rank above $k$.

The relation \eqref{eqn:incidence-plucker} states that
\[
  \min_{a\in T\setminus S}\bigl\{\theta(S+a)+\underline M(T-a)\bigr\}\ \text{vanishes tropically,}
  \qquad T\setminus S=\{x,y,z_1,\dots,z_{r-k}\}.
\]
We compute the summands $\underline M(T-a)$, which is $0$ if $T-a$ is a basis of $M$ and $\infty$
otherwise. The set $T-a$ has $r$ elements.
\begin{itemize}
\item For $a=y$: $T-y=(B-x)\cup\{x,z_1,\dots,z_{r-k}\}=B\cup\{z_1,\dots,z_{r-k}\}$, a basis of $M$
by choice; so $\underline M(T-y)=0$, and $\theta(S+y)=\theta(B-x+y)=\theta(B')$.
\item For $a=x$: $T-x=(B-x)\cup\{y,z_1,\dots,z_{r-k}\}=B'\cup\{z_1,\dots,z_{r-k}\}$ (using
$B'=B-x+y$). Since $\cl_M(B')=F=\cl_M(B)$ and $\{z_1,\dots,z_{r-k}\}$ extends $B$ to an $M$-basis,
it likewise extends $B'$ to an $M$-basis; so $T-x$ is a basis of $M$, $\underline M(T-x)=0$, and
$\theta(S+x)=\theta(B)$.
\item For $a=z_j$: $T-z_j=B\cup\{y\}\cup\{z_i:i\ne j\}$. Since $y\in F=\cl_M(B)$, this set has
$M$-rank $\le k+(r-k-1)=r-1<r$, hence is \emph{not} a basis of $M$; so $\underline M(T-z_j)=\infty$,
and the summand for $a=z_j$ is $\infty$.
\end{itemize}
Thus the only possibly-finite summands are those for $a=x$ and $a=y$, equal to $\theta(B)$ and
$\theta(B')$ respectively. For the tropical minimum to vanish (be attained at least twice, or be all
$\infty$), we need $\theta(B)=\theta(B')$. Hence $\theta(B)=\theta(B')$, completing the proof of
constancy. Well-definedness of $\widehat\theta$ and the identity $\theta(B)=\widehat\theta(\cl_M B)$
follow immediately.
\end{proof}

\Cref{prop:constancy} attaches to each $\theta\in\bV^{\,r-k}(M)$ its flat-coordinate
$\widehat\theta\in\cW_k(M)$, where $\cW_k(M)$ denotes the set of functions
$\cLL_k(M)\to\Rbar$ (not identically $\infty$) modulo global additive constant. Conversely, to a
function $w\colon\cLL_k(M)\to\Rbar$ we associate the function
\begin{equation}\label{eqn:thetaw}
  \theta_w\colon\binom{E}{k}\to\Rbar,\qquad
  \theta_w(B):=\begin{cases}w(\cl_M B), & \rk_M(B)=k,\\ \infty, & \rk_M(B)<k.\end{cases}
\end{equation}
A function $\theta\colon\binom{E}{k}\to\Rbar$ is \emph{flat-constant} (with respect to $M$) if
$\theta=\theta_w$ for some $w$. By \Cref{prop:constancy} every $\theta\in\bV^{\,r-k}(M)$ is
flat-constant and equals $\theta_{\widehat\theta}$ on its support, and $w\mapsto\theta_w$,
$\theta\mapsto\widehat\theta$ are mutually inverse between flat-constant functions and functions on
$\cLL_k(M)$ (with the support convention $\widehat{\theta_w}=w$ on $\operatorname{supp}w$). We write
\[
  \cW(M):=\bigsqcup_{0\le k\le r}\cW_k(M),\qquad
  \cW^{\Dr}_k(M):=\{\widehat\theta:\theta\in\bV^{\,r-k}(M)\}\subseteq\cW_k(M),
\]
and $\cW^{\Dr}(M):=\bigsqcup_k\cW^{\Dr}_k(M)$. By the previous paragraph and \Cref{fact:dressian},
the assignment $\theta\mapsto\widehat\theta$ is a poset isomorphism
\begin{equation}\label{eqn:identification}
  \Dr(M)\ \cong\ \bV(M)/\!\sim\ \xrightarrow{\ \cong\ }\ \cW^{\Dr}(M),
\end{equation}
the order on $\cW^{\Dr}(M)$ being transported from the quotient order. The membership problem ``which
$w\in\cW_k(M)$ lie in $\cW^{\Dr}_k(M)$'' is, by \eqref{eqn:thetaw}, exactly the requirement that
$\theta_w$ be a valuated quotient of $\underline M$; this is governed by the incidence relations
\eqref{eqn:incidence-plucker} with $\mu=\underline M$ and $m=r$, which we record once for reference:
\begin{equation}\label{eqn:incid-r}
  \min_{a\in T\setminus S}\bigl\{\theta_w(S+a)+\underline M(T-a)\bigr\}\ \text{vanishes tropically}
  \qquad\bigl(S\in\tbinom{E}{k-1},\ T\in\tbinom{E}{r+1}\bigr),
\end{equation}
together with the Pl\"ucker relations \eqref{eqn:plucker} for $\theta_w$.

\section{The candidate involution \texorpdfstring{$\Phi$}{Phi}}

\begin{definition}\label{def:Phi}
For $w\in\cW_k(M)$ define $w^{\perp}\in\cW_{r-k}(M)$ by
\begin{equation}\label{eqn:Phi-flat}
  w^{\perp}(F):=w(F^{\perp}),\qquad F\in\cLL_{r-k}(M),
\end{equation}
which is well-defined since $\perp\colon\cLL_{r-k}(M)\to\cLL_k(M)$ is a bijection (\Cref{lem:rank}).
This is an involution of the ambient space $\cW(M)$ exchanging $\cW_k(M)$ and $\cW_{r-k}(M)$. We
define the candidate map
\[
  \Phi\colon \cW(M)\to\cW(M),\qquad \Phi(w):=w^{\perp},
\]
and, abusing notation, also write $\Phi$ for the induced map on $\Dr(M)$ via \eqref{eqn:identification}
\emph{wherever it is defined}, i.e.\ on the subset of $\Dr(M)$ whose flat-coordinate $w$ has
$w^{\perp}\in\cW^{\Dr}(M)$.
\end{definition}

For $\Phi$ to descend to a self-map of $\Dr(M)$ we must know that $w^{\perp}$ is again the
flat-coordinate of a valuated quotient whenever $w$ is. This is the well-definedness statement, which
we now isolate.

\begin{claim}[Well-definedness, isolated]\label{conj:wd-invariance}
For every $0\le k\le r$ and every $w\in\cW^{\Dr}_k(M)$, one has $w^{\perp}\in\cW^{\Dr}_{r-k}(M)$.
Equivalently: if $\theta_w$ is a valuated quotient of $\underline M$ of rank $k$, then
$\theta_{w^{\perp}}$ is a valuated quotient of $\underline M$ of rank $r-k$.
\end{claim}

We are unable to prove \Cref{conj:wd-invariance} in full generality; see \Cref{rmk:honest} for a
precise account of why the natural reductions fail, and \Cref{rmk:numeric} for the cases in which we
have verified it. The remaining results of this document are conditional on
\Cref{conj:wd-invariance} where indicated, except for \Cref{prop:ext}, which is unconditional.

\begin{proposition}[Involution, conditional on \Cref{conj:wd-invariance}]\label{prop:inv}
Assume \Cref{conj:wd-invariance}. Then $\Phi$ is a well-defined self-map of $\Dr(M)$ and
$\Phi\circ\Phi=\operatorname{id}_{\Dr(M)}$.
\end{proposition}

\begin{proof}
By \Cref{conj:wd-invariance}, $w\mapsto w^{\perp}$ carries $\cW^{\Dr}(M)$ into itself; via
\eqref{eqn:identification} it therefore defines a self-map of $\Dr(M)$, which is $\Phi$.

For involutivity we work with honest representatives $w\colon\cLL_k(M)\to\Rbar$. The operation
\eqref{eqn:Phi-flat} is \emph{precomposition} with the index-set bijection $\perp$. Two facts about
precomposition give the result.

\emph{(i) It descends to scaling classes.} If $w'=w+c$ for a constant $c\in\R$ (with
$\infty+c=\infty$), then for every $F\in\cLL_{r-k}(M)$,
$(w')^{\perp}(F)=w'(F^{\perp})=w(F^{\perp})+c=w^{\perp}(F)+c$, so $(w+c)^{\perp}=w^{\perp}+c$. Thus
$[w]\mapsto[w^{\perp}]$ is well-defined on $\sim$-classes.

\emph{(ii) It is a strict involution on representatives.} For every $F\in\cLL_k(M)$, applying
\eqref{eqn:Phi-flat} twice and then $(F^{\perp})^{\perp}=F$ (\Cref{conv:perp}),
\[
  (w^{\perp})^{\perp}(F)=w^{\perp}(F^{\perp})=w\bigl((F^{\perp})^{\perp}\bigr)=w(F),
\]
so $(w^{\perp})^{\perp}=w$ as functions, with no additive shift.

Combining, for any class $[w]\in\cW^{\Dr}_k(M)$ we get
$(\Phi\circ\Phi)([w])=[(w^{\perp})^{\perp}]=[w]$. Under the poset isomorphism
\eqref{eqn:identification}, $\Phi\circ\Phi$ corresponds to the identity, so
$\Phi\circ\Phi=\operatorname{id}_{\Dr(M)}$.
\end{proof}

\begin{remark}[On order-reversal]\label{rmk:rev}
If \Cref{conj:wd-invariance} holds, then $\Phi$ is order-reversing as well. Indeed, order-reversal is
the assertion that for valuated quotients $\theta_2\onto\theta_1$ (ranks $k_2\ge k_1$, flat
coordinates $w_2,w_1$) one has $\theta_{w_1^{\perp}}\onto\theta_{w_2^{\perp}}$ (ranks $r-k_1\le
r-k_2$, the correct direction). The relation $\theta_2\onto\theta_1$ is the incidence system
\eqref{eqn:incidence-plucker} for the pair $(\theta_1,\theta_2)$, a system of relations indexed by
pairs $(S,T)$ with $|S|=k_1-1$, $|T|=k_2+1$. Precomposition of \emph{both} flat-coordinates with the
involutive bijection $\perp$ permutes the finite summands of this system among themselves (each set
$S+a$, $T-a$ is replaced by sets of complementary rank under $\perp$, and the two quotients trade
roles), so the transported system is the incidence system for $(\theta_{w_2^{\perp}},
\theta_{w_1^{\perp}})$. Making this matching precise requires the same control over the index sets as
\Cref{conj:wd-invariance} itself --- in fact well-definedness is the special case
$\theta_2=\underline M$, $k_2=r$ --- and we therefore do not claim order-reversal independently of
\Cref{conj:wd-invariance}. We have verified order-reversal computationally for $F_7$
(\Cref{rmk:numeric}).
\end{remark}

\section{The extension property (unconditional)}

\begin{proposition}[Extension of $\perp$]\label{prop:ext}
Under the inclusion $\cLL(M)^{\op}\into\Dr(M)$, $F\mapsto L_F:=\Trop(M/F\oplus U_{0,F})$, one has
$\Phi(L_F)=L_{F^{\perp}}$ for every flat $F$. In particular $\Phi$ is defined on the embedded copy of
$\cLL(M)^{\op}$ \emph{independently of \Cref{conj:wd-invariance}}, and there restricts to
$F\mapsto F^{\perp}$.
\end{proposition}

\begin{proof}
Fix a flat $F$ with $\rk_M(F)=s$ and set $k:=r-s$. The matroid $M/F\oplus U_{0,F}$ has rank $r-s=k$;
its bases are exactly the bases of the contraction $M/F$ (subsets of $E\setminus F$), the elements of
$F$ being loops. Its trivial valuation is $\{0,\infty\}$-valued, so by \Cref{prop:constancy} the
flat-coordinate $\widehat{L_F}\in\cW_{k}(M)$ is $\{0,\infty\}$-valued, with $\widehat{L_F}(K)=0$
exactly for those $K\in\cLL_{k}(M)$ that arise as $\cl_M(B)$ for a basis $B$ of $M/F$. A $k$-subset
$B\subseteq E\setminus F$ is a basis of $M/F$ iff $B$ is independent in $M$ and $\cl_M(B\cup F)=E$;
for such $B$, $K:=\cl_M(B)$ is a rank-$k$ flat with
\begin{equation}\label{eqn:complement}
  K\vee F=\hat1\quad\text{and}\quad K\wedge F=\hat0 ,
\end{equation}
i.e.\ $K$ is a (modular) complement of $F$ of rank $k=r-s$; conversely every rank-$k$ flat $K$
satisfying \eqref{eqn:complement} is the $M$-closure of a basis of $M/F$ (take a basis of $M|_K$: it
is disjoint from $F$ by $K\wedge F=\hat0$, independent in $M/F$, and spanning by $K\vee F=\hat1$).
Hence
\begin{equation}\label{eqn:LF-coord}
  \widehat{L_F}(K)=\begin{cases}0,& \rk_M(K)=r-s,\ K\vee F=\hat1,\ K\wedge F=\hat0,\\
  \infty,&\text{otherwise.}\end{cases}
\end{equation}

Now apply \eqref{eqn:Phi-flat}. For $K'\in\cLL_{s}(M)$,
$\widehat{L_F}^{\,\perp}(K')=\widehat{L_F}(K'^{\perp})$, and by \Cref{lem:rank} $K'^{\perp}$ ranges
over $\cLL_{r-s}(M)$ as $K'$ ranges over $\cLL_{s}(M)$. By \eqref{eqn:LF-coord},
$\widehat{L_F}(K'^{\perp})=0$ iff $K'^{\perp}\vee F=\hat1$ and $K'^{\perp}\wedge F=\hat0$. Applying
the anti-automorphism $\perp$ (which exchanges $\vee\leftrightarrow\wedge$ and
$\hat0\leftrightarrow\hat1$, by \Cref{lem:rank}) and using $(K'^{\perp})^{\perp}=K'$:
\[
  K'^{\perp}\vee F=\hat1\ \Longleftrightarrow\ K'\wedge F^{\perp}=\hat0,\qquad
  K'^{\perp}\wedge F=\hat0\ \Longleftrightarrow\ K'\vee F^{\perp}=\hat1 .
\]
Therefore $\widehat{L_F}^{\,\perp}(K')=0$ iff $K'$ is a rank-$s$ flat with $K'\vee F^{\perp}=\hat1$
and $K'\wedge F^{\perp}=\hat0$; since $\rk_M(F^{\perp})=r-s$ (\Cref{lem:rank}), this says exactly that
$K'$ is a rank-$s$ modular complement of $F^{\perp}$. Comparing with \eqref{eqn:LF-coord} applied to
the flat $F^{\perp}$ (rank $r-s$, so $r-\rk_M(F^{\perp})=s$),
\[
  \widehat{L_{F^{\perp}}}(K')=\begin{cases}0,&\rk_M(K')=s,\ K'\vee F^{\perp}=\hat1,\ K'\wedge
  F^{\perp}=\hat0,\\ \infty,&\text{otherwise,}\end{cases}
\]
we obtain $\widehat{L_F}^{\,\perp}=\widehat{L_{F^{\perp}}}$ in $\cW_{s}(M)$. This identity is at the
level of \emph{honest flat-coordinate functions}, so no appeal to \Cref{conj:wd-invariance} is
needed; both $\widehat{L_F}^{\,\perp}$ and $\widehat{L_{F^{\perp}}}$ are manifestly flat-coordinates
of valuated quotients (the latter is $\widehat{L_{F^{\perp}}}\in\cW^{\Dr}(M)$ by construction). By
\eqref{eqn:identification} this gives $\Phi(L_F)=L_{F^{\perp}}$.
\end{proof}

\section{Honest status and computational evidence}

\begin{remark}[What is proved, and what is the precise gap]\label{rmk:honest}
The following are proved here in full, using only the foundational facts \Cref{fact:dressian} and
\Cref{fact:dual} (whose hypotheses we verify in the present setting):
\begin{enumerate}
\item the rank identity and the lattice anti-automorphism property of $\perp$ (\Cref{lem:rank});
\item the canonical flat-coordinate $\widehat\theta$ of every element of $\Dr(M)$, via
\emph{constancy on flats} (\Cref{prop:constancy}); the proof uses crucially that $\theta$ is a
quotient of the \emph{trivial} valuation $\underline M$ (the auxiliary elements $z_j$ in
\eqref{eqn:incid-r} kill the off-flat terms);
\item the explicit candidate $\Phi$ as precomposition $w\mapsto w\circ\perp$ on flat-coordinates
(\Cref{def:Phi}), which is type-correct because source and target are functions on flats of
complementary rank, identified by $\perp$;
\item the \emph{extension property} $\Phi(L_F)=L_{F^{\perp}}$ (\Cref{prop:ext}), proved
unconditionally at the level of flat-coordinate functions from \eqref{eqn:LF-coord} and the lattice
anti-automorphism property of $\perp$;
\item conditionally on the single isolated \Cref{conj:wd-invariance}, that $\Phi$ is a well-defined
self-map of $\Dr(M)$ and an involution (\Cref{prop:inv}); and (\Cref{rmk:rev}) order-reversing.
\end{enumerate}
The entire remaining gap is \Cref{conj:wd-invariance}: the property ``$\theta_w$ is a valuated
quotient of $\underline M$'' is invariant under $w\mapsto w\circ\perp$. This is a statement about the
incidence system \eqref{eqn:incid-r} (and the Pl\"ucker relations \eqref{eqn:plucker} for $\theta_w$);
its difficulty is genuine, as explained in \Cref{rmk:false-reduction}.
\end{remark}

\begin{remark}[Why the natural ``essential relations'' reduction fails]\label{rmk:false-reduction}
A tempting route to \Cref{conj:wd-invariance} is to find a finite, manifestly $\perp$-invariant
system of relations on $w\in\cW_k(M)$ that characterises $\cW^{\Dr}_k(M)$, modelled on the
three-term tropical Pl\"ucker relations. The most natural candidate is the system of
\emph{modular three-term relations}: for each interval $[A,Z]$ of $\cLL(M)$ with $\rk_M(A)=k-1$ and
$\rk_M(Z)=k+1$,
\[
  \min_{F\in\cLL_k(M),\,A<F<Z} w(F)\ \text{vanishes tropically.}
\]
This system \emph{is} $\perp$-invariant (it is carried to the corresponding rank-$(r-k)$ system by
$\perp$, since $\perp$ maps each interval $[A,Z]$ to the interval $[Z^{\perp},A^{\perp}]$ and the
rank-$k$ atoms strictly between $A,Z$ to the rank-$(r-k)$ atoms strictly between $Z^{\perp},A^{\perp}$).
\emph{However, it does not characterise $\cW^{\Dr}_k(M)$.} Two concrete failures, which we have
verified by exhaustive computation:
\begin{itemize}
\item For $M=U_{2,5}$ and $k=1$: the function $w$ on the five rank-$1$ flats (points) with values
$(1,0,0,0,1)$ satisfies the modular three-term relations (the unique interval $[\hat0,\hat1]$ has all
five points as atoms, and $\min=0$ is attained three times), yet $\theta_w$ is \emph{not} a valuated
matroid: the Pl\"ucker relation \eqref{eqn:plucker} at the pair $T=\{p_1,p_2\}$ with values $1,0$ has
minimum attained once. Thus the modular relations are strictly weaker than valuated-matroid-hood.
\item For $M=U_{3,5}$ and $k=2$ there are flat-constant functions satisfying both the modular
relations and the matroid-quotient support condition that nonetheless violate the incidence relations
\eqref{eqn:incid-r}; we exhibited $>2000$ such functions among the $3^{10}$ tested.
\end{itemize}
The root cause is structural: the incidence relations \eqref{eqn:incid-r} are indexed by pairs
$(S,T)$ with $|T|=r+1$, and the multiset of flats $\{\cl_M(S+a):a\in T\setminus S,\ T-a\text{ a
basis}\}$ that co-occur in a single relation is \emph{finer} than ``all rank-$k$ flats in a
length-two interval.'' Consequently the relations cutting out $\cW^{\Dr}_k(M)$ do not, in general,
form a $\perp$-invariant system of \emph{local} three-term type, and a different (global) argument is
needed. This is exactly the obstruction that \Cref{conj:wd-invariance} encapsulates and that we
believe constitutes the technical content of the announced reference.
\end{remark}

\begin{remark}[Computational verification of the full Theorem in the base case]\label{rmk:numeric}
We have verified the complete statement of the Theorem --- including the otherwise-unproved
\Cref{conj:wd-invariance} --- for the Fano matroid $M=F_7=\mathrm{PG}(2,2)$ (rank $3$, $7$ points), a
modular matroid, equipped with the polarity $\perp$ induced by the standard symmetric bilinear form
on $\mathbb F_2^3$ (so a point $p$ maps to the line $\{q:p\cdot q=0\}$, and conversely). Concretely,
over all flat-coordinates with small integer (and $\infty$) values we checked, by direct evaluation
of \eqref{eqn:plucker} and \eqref{eqn:incid-r}:
\begin{itemize}
\item $\perp$ is an order-reversing involution of $\cLL(F_7)$ with $\rk(F)+\rk(F^{\perp})=3$
(confirming \Cref{lem:rank});
\item \emph{constancy on flats} (\Cref{prop:constancy}) holds for every valuated quotient sampled;
\item \emph{well-definedness} (\Cref{conj:wd-invariance}): for every valuated quotient $\theta_w$ of
$\underline{F_7}$ (of each rank $k\in\{1,2\}$), the function $\theta_{w^{\perp}}$ is again a valuated
quotient of $\underline{F_7}$ of rank $3-k$ --- $0$ failures over $65$ rank-$2$ and $171$ rank-$1$
quotients with values in $\{0,1,(2),\infty\}$;
\item $\Phi\circ\Phi=\operatorname{id}$ on the nose (\Cref{prop:inv});
\item \emph{order-reversal} (\Cref{rmk:rev}): over all $837$ quotient-pairs $\theta_2\onto\theta_1$
found, $\theta_{w_1^{\perp}}\onto\theta_{w_2^{\perp}}$ held with $0$ failures;
\item \emph{extension} (\Cref{prop:ext}): $\Phi(L_F)=L_{F^{\perp}}$ for every flat $F$ of $F_7$.
\end{itemize}
We additionally confirmed, on $U_{2,5},U_{2,6},U_{3,5},M(K_4)$, the constancy statement
(\Cref{prop:constancy}) and the two counterexamples of \Cref{rmk:false-reduction}. These computations
support the truth of the Theorem and of \Cref{conj:wd-invariance}, while making clear that the latter
is not reducible to the local relations and remains the open analytic core.
\end{remark}

\begin{remark}[Boolean base case]\label{rmk:Unn}
When $M=U_{n,n}$, every subset of $[n]$ is a flat, $\cl_M(B)=B$, $\cLL_k(M)=\binom{[n]}{k}$, and the
flat-coordinate maps are the identity, so $\bV^{\,n-k}(U_{n,n})/\!\sim$ is the full set of rank-$k$
valuated matroids on $[n]$ up to scaling, i.e.\ $\Dr(U_{n,n})$ is the classical Dressian of $\PT{n}$.
Taking $\perp=$ set complementation $F\mapsto E\setminus F$ (the unique self-complementing
anti-automorphism of the Boolean lattice), \eqref{eqn:Phi-flat} reads $w^{\perp}(B)=w(E\setminus B)$,
which is precisely the Dress--Wenzel valuated dual $\theta\mapsto\theta^{*}$ of \Cref{fact:dual}. In
this case \Cref{conj:wd-invariance}, involutivity, and order-reversal all hold by \Cref{fact:dual},
so the Theorem is established \emph{unconditionally} for $U_{n,n}$, and $\Phi$ agrees on the nose with
classical valuated matroid duality, with no extra sign or rescaling. This is the motivating instance,
and it shows that the general $\Phi$ of \Cref{def:Phi} is the correct relative analogue.
\end{remark}

\begin{proof}[Proof of the Theorem (conditional)]
By \eqref{eqn:identification}, $\Dr(M)$ is identified as a poset with $\cW^{\Dr}(M)$. Assuming
\Cref{conj:wd-invariance}, the map $\Phi$ of \Cref{def:Phi} is a well-defined self-map of $\Dr(M)$,
an involution (\Cref{prop:inv}), and order-reversing (\Cref{rmk:rev}); and unconditionally it
extends $F\mapsto F^{\perp}$ on the embedded $\cLL(M)^{\op}$ (\Cref{prop:ext}). For $M=U_{n,n}$,
\Cref{conj:wd-invariance} holds by valuated duality (\Cref{rmk:Unn}), giving the Theorem
unconditionally there; for general modular $M$ admitting a polarity, the Theorem is verified
computationally in the base case $F_7$ (\Cref{rmk:numeric}), and is reduced to the single isolated
\Cref{conj:wd-invariance}.
\end{proof}

\end{document}
